\documentclass[a4paper,10pt]{ctexart}
\usepackage{amsmath,amssymb}
\usepackage[margin=2cm]{geometry}
\usepackage{tikz}
\usetikzlibrary{calc,patterns,decorations.pathmorphing,arrows.meta}
\usepackage{graphicx}
\usepackage{float}
\begin{document}
\noindent 1. 【2024·山东】如图所示，质量均为$m$的甲、乙两同学，分别坐在水平放置的轻木板上，木板通过一根原长为$l$的轻质弹性绳连接，连接点等高且间距为$d$（$d < l$）。两木板与地面间动摩擦因数均为$\mu$，弹性绳劲度系数为$k$，被拉伸时弹性势能$E_p = \frac{1}{2}kx^2$（$x$为绳的伸长量）。现用水平力$F$缓慢拉动乙所坐木板，直至甲所坐木板刚要离开原位置，此过程中两人与所坐木板保持相对静止，$k$保持不变，最大静摩擦力等于滑动摩擦力，重力加速度大小为$g$，则$F$所做的功等于（\quad）
\begin{center}
	
	
	\begin{tikzpicture}[>=Stealth,scale=1]
		% 公共参数
		\def\blocksize{0.6} % 方块边长
		\def\springwidth{0.8} % 弹簧宽度
		\def\groundlength{7} % 地面长度
		\def\blockmarg{0.1} % 方块与地面间距
		
		% ---------------- 第一幅图 ----------------
		\begin{scope}[yshift=0cm]
			% 地面
			\draw[pattern=north east lines] (0,0) rectangle (\groundlength, -0.1);
			\draw (0,0) -- (\groundlength,0);
			% 甲、乙方块
			\draw (0.5,0) rectangle (0.5+\blocksize,\blocksize);
			\node at (0.5+\blocksize/2,\blocksize+0.2) {$m$};
			\node at (0.5+\blocksize/2,\blocksize/2) {甲};
			
			\draw (0.5+3,0) rectangle (0.5+3+\blocksize,\blocksize);
			\node at (0.5+3+\blocksize/2,\blocksize+0.2) {$m$};
			\node at (0.5+3+\blocksize/2,\blocksize/2) {乙};
			% 弹簧
			\draw[decorate,decoration={coil,amplitude=\springwidth,segment length=3pt}] 
			(0.5+\blocksize,\blocksize/2) -- (0.5+3,\blocksize/2) node[midway,above] {$k$};
			% 距离标注d
			\draw[|<->|] (0.5+\blocksize+0.1,\blocksize+0.3) -- (0.5+3-0.1,\blocksize+0.3) node[midway,above] {$d$};
			% 拉力F
			\draw[->] (0.5+3+\blocksize+0.2,\blocksize/2) -- (0.5+3+\blocksize+1,\blocksize/2) node[right] {$F$};
			% 摩擦系数μ
			\node at (\groundlength-0.5,0.2) {$\mu$};
		\end{scope}
		
		% ---------------- 第二幅图（原长） ----------------
		\begin{scope}[yshift=-2cm]
			% 甲、乙方块
			\draw (0.5,0) rectangle (0.5+\blocksize,\blocksize);
			\node at (0.5+\blocksize/2,\blocksize/2) {甲};
			
			\draw (0.5+4,0) rectangle (0.5+4+\blocksize,\blocksize);
			\node at (0.5+4+\blocksize/2,\blocksize/2) {乙};
			% 弹簧原长线
			\draw[decorate,decoration={coil,amplitude=\springwidth,segment length=3pt}]	(0.5+\blocksize,\blocksize/2) -- (0.5+4,\blocksize/2);
			% 标注原长l
			\draw[|<->|] (0.5+\blocksize+0.1,\blocksize+0.3) -- (0.5+4-0.1,\blocksize+0.3) node[midway,above] {$l$};
			\node at (0.5+\blocksize/2+2,\blocksize/2-0.75) {原长};
			% 拉力F
			\draw[->] (0.5+4+\blocksize+0.2,\blocksize/2) -- (0.5+4+\blocksize+1,\blocksize/2) node[right] {$F$};
			% 地面
			\draw[pattern=north east lines] (0,0) rectangle (\groundlength, -0.1);
		\end{scope}
		
		% ---------------- 第三幅图（刚要离开） ----------------
		\begin{scope}[yshift=-4cm]
			% 甲、乙方块
			\draw (0.5,0) rectangle (0.5+\blocksize,\blocksize);
			\node at (0.5+\blocksize/2,\blocksize/2) {甲};
			\node at (0.5+\blocksize/2,-0.5) {刚要离开};
			
			\draw (0.5+5,0) rectangle (0.5+5+\blocksize,\blocksize);
			\node at (0.5+5+\blocksize/2,\blocksize/2) {乙};
			% 弹簧线
			\draw[decorate,decoration={coil,amplitude=\springwidth,segment length=3pt}] (0.5+\blocksize,\blocksize/2) -- (0.5+5,\blocksize/2);
			% 弹性势能公式
			\node at (3,0.75) {$E_p = \dfrac{1}{2}kx^2$};
			% 拉力F
			\draw[->] (0.5+5+\blocksize+0.2,\blocksize/2) -- (0.5+5+\blocksize+1,\blocksize/2) node[right] {$F$};
			% 地面
			\draw[pattern=north east lines] (0,0) rectangle (\groundlength, -0.1);
		\end{scope}
	\end{tikzpicture}
\end{center}	

\noindent \textbf{解析：}

设弹性绳的拉力为$T$。

开始时两连接点间距$d<l$，弹性绳处于松弛状态，因此在乙被缓慢拉动的前一段过程中，甲始终不受绳的拉力，乙只需克服摩擦力前进。

当甲所坐木板刚要离开原位置时，甲受到的静摩擦力达到最大值，所以有
\[
T=\mu mg.
\]

由胡克定律可得此时弹性绳的伸长量
\[
x=\frac{T}{k}=\frac{\mu mg}{k}.
\]

乙所坐木板的总位移分两段：

第一段：把两连接点间距从$d$拉到$l$，位移为
\[
s_1=l-d.
\]
此过程中弹性绳没有拉力，拉力$F$只需平衡乙受到的摩擦力，故
\[
F_1=\mu mg.
\]
这一段外力做功
\[
W_1=F_1s_1=\mu mg(l-d).
\]

第二段：弹性绳由原长继续伸长$x$。由于过程缓慢，乙始终近似平衡，故任意时刻外力
\[
F=\mu mg+T.
\]
其中$T$随伸长量从$0$增大到$\mu mg$。

因此第二段外力做功等于克服乙的摩擦力做功与增加弹性势能之和：
\[
W_2=\mu mg\cdot x+\frac{1}{2}kx^2.
\]

代入$x=\dfrac{\mu mg}{k}$，得
\[
W_2=\mu mg\cdot \frac{\mu mg}{k}+\frac{1}{2}k\left(\frac{\mu mg}{k}\right)^2
=\frac{3\mu^2m^2g^2}{2k}.
\]

所以$F$所做的总功为
\[
W=W_1+W_2=\mu mg(l-d)+\frac{3\mu^2m^2g^2}{2k}.
\]

\noindent \textbf{答案：}
\[
\boxed{W=\mu mg(l-d)+\frac{3\mu^2m^2g^2}{2k}}
\]


\noindent 2. 【2024·贵州】质量为1kg的物块静置于光滑水平地面上，设物块静止时的位置为$x$轴零点。现给物块施加一沿$x$轴正方向的水平力$F$，其大小随位置$x$变化的关系如图所示，则物块运动到$x=3\,\text{m}$处，$F$做功的瞬时功率为（\quad）

A. $8\,\text{W}$ \quad B. $16\,\text{W}$ \quad C. $24\,\text{W}$ \quad D. $36\,\text{W}$	

\begin{center}
	\begin{tikzpicture}[>=stealth,scale=1.2]
		% 坐标轴
		\draw[->] (0,0) -- (4,0) node[right] {$x/\text{m}$};
		\draw[->] (0,0) -- (0,4) node[above] {$F/\text{N}$};
		% 刻度
		\foreach \x in {0,1,2,3}
		\draw (\x,0.1) -- (\x,-0.1) node[below] {$\x$};
		\foreach \y in {1,2,3}
		\draw (0.1,\y) -- (-0.1,\y) node[left] {$\y$};
		% 第一段：0~2m，F=3N
		\draw[thick] (0,3) -- (2,3);
		\draw[dashed] (2,0) -- (2,3);
		% 第二段：2~3m，F=2N
		\draw[thick] (2,2) -- (3,2);
		\draw[dashed] (2,2) -- (0,2);
		\draw[dashed] (3,0) -- (3,2);
	\end{tikzpicture}
\end{center}	

\noindent 3.

\noindent 如图甲所示，固定粗糙斜面的倾角为$37^\circ$，与斜面平行的轻弹簧下端固定在$C$处，上端连接质量为$1\mathrm{kg}$的小滑块（视为质点），$BC$为弹簧的原长。现将滑块从$A$处由静止释放，在滑块从释放至第一次到达最低点的过程中，其加速度$a$随弹簧的形变量$x$的变化规律如图乙所示（取沿斜面向下为加速度的正方向），取$\sin37^\circ=0.6$，$\cos37^\circ=0.8$，$g=10\mathrm{m/s^2}$。根据上述过程，试求出下列各量：

\begin{enumerate}
	\item 滑块与斜面间的动摩擦因数$\mu$；
	\item 当$x=0$时滑块的动能$E_\text{k}$；
	\item 弹簧的最大弹性势能$E_\text{p}$。
\end{enumerate}

% 题目配图（图甲：斜面+弹簧+滑块）
\begin{figure}[H]
\centering
\begin{tikzpicture}[>=Stealth,scale=1]
	% 斜面与地面
	\coordinate (O) at (0,0);
	\coordinate (D) at (4,0);
	\coordinate (S) at (4,3);
	\draw[thick] (O) -- (D);
	\draw[thick] (O) -- (S);

	% 倾角标注
	\draw (1.2,0) arc[start angle=0,end angle=37,radius=1.2];
	\node at (1.55,0.45) {$37^\circ$};

	% 斜面上的关键点
	\coordinate (C) at (0,0);
	\coordinate (B) at (3,2.25);
	\coordinate (A) at (4,3);

	% 弹簧与滑块
	\draw[decorate,decoration={coil,segment length=4pt,amplitude=2.4pt}] ($(C)+(0,0.2)$) -- ($(A)+(0,0.2)$);
	\begin{scope}[shift={(A)},rotate=37]
		\draw[fill=white] (-0.28,0) rectangle (0.28,0.42);
	\end{scope}

	% 点与文字标注
	\fill (A) circle (1.2pt);
	\fill (B) circle (1.2pt);
	\fill (C) circle (1.2pt);
	\node[above left,rotate=37] at ($(A)+(0.2,0.1)$) {$A$};
	\node[above,rotate=37] at ($(B)+(-0.1,0.3)$) {$B$};
	\node[below left] at (C) {$C$};
	\node[right] at ($(A)+(0.35,0.1)$) {滑块};

	% BC 为原长
	\draw[|<->|] ($(B)+(-0.32,0.43)$) -- ($(C)+(-0.32,0.43)$)
		node[midway,sloped,above] {$BC$为原长};

	\node at (2.4,-0.7) {图甲};
\end{tikzpicture}	
	% 图乙：a-x图像
	\begin{tikzpicture}[x=10cm,y=1cm,scale=0.5]
		\coordinate (A) at (0.2,4);
		\coordinate (B) at (0,2);
		\coordinate (C) at (0.2,0);
		\coordinate (D) at (0.6,-4);
		% 坐标轴
		\draw[->] (0,-4.5) -- (0,4.5) node[above] {$a\ (\mathrm{m\cdot s^{-2}})$};
		\draw[->] (-0.1,0) -- (0.7,0) node[right] {$x\ (\mathrm{m})$};
		% 刻度
		\foreach \y in {-4,-3,-2,-1,1,2,3,4} {
			\draw (0.02,\y) -- (-0.02,\y) node[left] {$\y$};
		}
		\foreach \x in {0.2,0.4,0.6} {
			\draw (\x,0.05) -- (\x,-0.05) node[below] {$\x$};
		}
		\draw(A)--(B)--(C)--(D);
		\node at (0.2,-5) {图乙};
\end{tikzpicture}
	\caption{图甲：斜面装置；图乙：$a-x$变化规律}

\end{figure}

\noindent \textbf{解析：}

设沿斜面向下为正方向。由图乙可知，当$x=0$时，滑块恰好运动到弹簧原长位置，此时加速度为
\[
a=2\,\mathrm{m/s^2}.
\]

此时弹簧弹力为零，由牛顿第二定律得
\[
mg\sin37^\circ-\mu mg\cos37^\circ=ma.
\]

代入$m=1\,\mathrm{kg}$，$g=10\,\mathrm{m/s^2}$，$\sin37^\circ=0.6$，$\cos37^\circ=0.8$，$a=2\,\mathrm{m/s^2}$，得
\[
10\times 0.6-\mu\cdot 10\times 0.8=2,
\]
\[
6-8\mu=2,
\]
\[
\mu=0.5.
\]

再看$x>0$时图像，可得直线经过$(0,2)$与$(0.2,0)$，故
\[
a=2-10x.
\]

而在弹簧被压缩时，有
\[
a=g\sin37^\circ-\mu g\cos37^\circ-\frac{kx}{m}=2-kx.
\]

与图像表达式比较得
\[
k=10\,\mathrm{N/m}.
\]

滑块由$A$运动到$B$（即$x$由$0.2\,\mathrm{m}$减小到$0$）的过程中，初速度为零，设到达$x=0$时动能为$E_\text{k}$。

由动能定理得
\[
E_\text{k}=mg\sin37^\circ\cdot 0.2-\mu mg\cos37^\circ\cdot 0.2+\frac{1}{2}k(0.2)^2.
\]

代入数据得
\[
E_\text{k}=6\times 0.2-4\times 0.2+\frac{1}{2}\times 10\times (0.2)^2
=1.2-0.8+0.2=0.6\,\mathrm{J}.
\]

设滑块第一次到达最低点时，弹簧压缩量为$x_\text{m}$，此时速度再次为零。

从$x=0$运动到最低点的过程中，由动能定理得
\[
0-E_\text{k}=mg\sin37^\circ\cdot x_\text{m}-\mu mg\cos37^\circ\cdot x_\text{m}-\frac{1}{2}kx_\text{m}^2.
\]

即
\[
-0.6=2x_\text{m}-5x_\text{m}^2.
\]

整理得
\[
5x_\text{m}^2-2x_\text{m}-0.6=0.
\]

解得
\[
x_\text{m}=0.6\,\mathrm{m}
\]
（负值舍去）。

故弹簧的最大弹性势能为
\[
E_\text{p}=\frac{1}{2}kx_\text{m}^2=\frac{1}{2}\times 10\times (0.6)^2=1.8\,\mathrm{J}.
\]

\noindent \textbf{答案：}
\[
\boxed{\mu=0.5}
\]
\[
\boxed{E_\text{k}=0.6\,\mathrm{J}}
\]
\[
\boxed{E_\text{p}=1.8\,\mathrm{J}}
\]

\newpage
\textbf{解答（高中方法，无积分）}
\begin{enumerate}
	\item 
\textbf{求滑块与斜面间的动摩擦因数$\mu$}

当$x=0$时，弹簧处于原长，弹力为$0$。
对滑块沿斜面方向由牛顿第二定律：
\[
mg\sin37^\circ - \mu mg\cos37^\circ = ma_0
\]
代入$m=1\mathrm{kg}$，$g=10\mathrm{m/s^2}$，$\sin37^\circ=0.6$，$\cos37^\circ=0.8$，$a_0=2\mathrm{m/s^2}$：
\begin{align*}
	1\times10\times0.6 - \mu\times1\times10\times0.8 &= 1\times2 \\
	6 - 8\mu &= 2 \\
	\mu &= 0.5
\end{align*}


\item 
\textbf{当$x=0$时滑块的动能$E_\text{k}$}

根据\textbf{动能定理}，合外力做功等于动能变化：$W_\text{合} = \Delta E_\text{k}$。
合外力$F_\text{合}=ma$，在$a-x$图像中，合外力做功等于$m$乘以$a-x$图线与$x$轴围成的面积（高中面积法）。

从$A$（$x=0.2\mathrm{m}$）到$B$（$x=0$），$a-x$图为梯形：
上底$a_1=2\mathrm{m/s^2}$，下底$a_2=4\mathrm{m/s^2}$，高$\Delta x=0.2\mathrm{m}$
\[
S = \frac{(a_1+a_2)}{2} \cdot \Delta x = \frac{(2+4)}{2} \times 0.2 = 0.6\ \mathrm{m\cdot s^{-2}}
\]
合外力做功：
$
W_\text{合} = mS = 1\times0.6 = 0.6\ \mathrm{J}
$
因滑块从静止释放，初动能为$0$，故$E_\text{k} = W_\text{合} = 0.6\mathrm{J}$。

\item 
\textbf{求弹簧的最大弹性势能$E_\text{p}$}

当滑块第一次到达最低点时，速度为$0$，弹簧形变量最大，设为$x_m$。
由图乙中$x>0$部分知，滑块压缩弹簧时的加速度满足
\[
a=2-10x.
\]

从$x=0$到第一次到达最低点，由动能定理可得合外力做功等于动能变化。
在$a-x$图像中，这对应于图线与$x$轴围成的有向面积乘以质量。

其中，$x=0$到$x=0.2$之间图线在$x$轴上方，围成一个三角形，面积为
\[
S_1=\frac12\times 0.2\times 2=0.2.
\]

$x=0.2$到$x=x_m$之间图线在$x$轴下方，围成一个三角形，面积为
\[
S_2=\frac12\times (x_m-0.2)\times (10x_m-2).
\]

因为此过程末动能为$0$，初动能为$E_\text{k}=0.6\mathrm{J}$，所以合外力做功为
\[
W_\text{合}=0-E_\text{k}=-0.6\mathrm{J}.
\]

又$m=1\mathrm{kg}$，故有
\[
S_1-S_2=-0.6.
\]

即
\[
0.2-\frac12(x_m-0.2)(10x_m-2)=-0.6.
\]

化简得
\[
5x_m^2-2x_m-0.6=0.
\]

整理得
\[
x_m=0.6\mathrm{m}
\]
（负值舍去）。

故弹簧的最大弹性势能为
\[
E_\text{p}=\frac12 kx_m^2=\frac12\times10\times(0.6)^2=1.8\mathrm{J}.
\]
\end{enumerate}
\noindent
\textbf{最终答案：}
(1) $\boxed{\mu=0.5}$；(2) $\boxed{E_\text{k}=0.6\mathrm{J}}$；(3) $\boxed{E_\text{p}=1.8\mathrm{J}}$
\end{document}